\section{Introduction}
\label{sec:intro}

Who bears the cost of a national resource plan, and when? Countries set their climate,
land, energy and water strategies with planning models that identify the least-cost way
to meet future needs within physical and policy limits. A plan can be technically sound
and still have very different consequences for households, industries and public
finances. Lower operating costs may require more investment today. Cleaner production
can improve health. Higher electricity costs can reduce purchasing power and industrial
activity, and change the very demand the plan was designed to meet.

Planning models answer what to build and how to operate it. They contain no households,
firms or government budget, so they cannot show who pays for a plan or who benefits from
it. Economic models that do contain these usually represent energy only in aggregate,
without the technologies, infrastructure and physical limits that planners work with.
The economic consequences of a resource plan therefore fall between two well-developed
modelling traditions.

OG--CLEWS connects the two. It links the Climate, Land, Energy and Water Systems model
(CLEWS) \citep{howells_clews_2013}, built on the OSeMOSYS optimiser
\citep{howells_osemosys_2011}, with OG-Core \citep{debacker_evans_ogcore}, an
overlapping-generations model of the economy. In OG-Core, households differ by age and
lifetime income, work, save for retirement and respond to prices and taxes; firms in
several industries employ workers and capital; and a government collects taxes, pays
transfers and pensions, invests in infrastructure and manages its debt. Linking the two
models allows the consequences of a resource plan to be followed through production,
household budgets and public finances, over time and across generations.

The link works in both directions. The resource plan changes costs, investment and
pollution, and these change the economy. The economy's response changes the demand for
electricity and the return on capital, which the resource plan should take into account.
Repeating the exchange allows the plan and the economy to become consistent with each
other.

Other approaches connect resource models with the economy in different ways.
Energy--economy models such as MARKAL--MACRO and MESSAGE--MACRO link energy costs with
growth and demand \citep{manne_wene_1992,messner_schrattenholzer_2000}, and computable
general equilibrium models add detail on industries and households
\citep{bohringer_rutherford_2008,rausch_2011}. More recent frameworks connect the economy
with environmental accounts and ecosystem services \citep{banerjee_ieem_2020}, or trace
employment through suppliers and household spending \citep{bararzadeh_clews_io_2026}.
Overlapping-generations models have been used to study the generational effects of
climate policy \citep{kotlikoff_2021}. OG--CLEWS brings that life-cycle and fiscal
perspective to a national resource plan with explicit technologies, infrastructure and
physical limits. We illustrate the framework with a Philippine energy-policy scenario.
